\chapter{Conclusion and Learning Path}
\label{chap:conclusion}

\section{The Integrated Mental Model}

The essential outcome of this handbook is a connected engineering worldview. Embedded AI begins with physical sensing and ends with physical action. Between those points, measurements are represented, aligned, fused, interpreted, propagated through state and prediction, converted into decisions, and executed by actuators. Resource and safety constraints cross every stage.

The main conceptual distinctions are:

\begin{itemize}
    \item a \textbf{sensor} is a physical device; a \textbf{modality} is a type of information;
    \item a \textbf{measurement} is an observation; a \textbf{state} is an inferred physical condition;
    \item a \textbf{representation} determines the structure available to the algorithm;
    \item \textbf{fusion} can occur at data, feature, decision, state, or runtime-management level;
    \item \textbf{tight coupling} shares measurements and residuals inside one estimator;
    \item \textbf{attention} performs learned input-dependent weighting but is not automatically uncertainty;
    \item \textbf{event-driven} processing and \textbf{spiking} processing are related but distinct;
    \item \textbf{model efficiency} is not the same as \textbf{system efficiency};
    \item \textbf{open-loop accuracy} does not guarantee \textbf{closed-loop safety}.
\end{itemize}

\section{A Practical Sequence for Further Study}

A productive learning sequence is:

\begin{enumerate}[label=\arabic*.]
    \item \textbf{Sensing physics:} understand what camera, event, LiDAR, radar, and IMU actually measure.
    \item \textbf{Representation and alignment:} study timestamps, coordinate frames, calibration, projection, and data association.
    \item \textbf{Fusion and estimation:} learn probabilistic fusion, Kalman filtering, factor graphs, and multimodal feature fusion.
    \item \textbf{AI architectures:} connect CNN, GNN, transformer, and SNN structures to their data and hardware.
    \item \textbf{Embedded systems:} measure latency, memory, energy, dataflow, and real-time scheduling.
    \item \textbf{Control and robotics:} understand how estimates become plans and actuator commands.
    \item \textbf{Safety and validation:} test failure modes, uncertainty, fallback, and closed-loop behavior.
\end{enumerate}

This sequence should be iterative rather than rigid. A paper reveals a gap in the mental model; the foundational concept is studied; the paper is reread with greater depth.

\section{The Week 7 Papers Revisited}

The six papers now form a connected set:

\begin{itemize}
    \item W7-1 aligns event sensing with asynchronous computation.
    \item W7-2 manages sensors and models as runtime energy resources.
    \item W7-3 redesigns multimodal attention for compute-in-memory constraints.
    \item W7-4 learns whether modality data should be transmitted at all.
    \item W7-5 uses tight coupling to preserve state observability.
    \item W7-6 connects feature attention to closed-loop driving behavior.
\end{itemize}

Their shared lesson is that intelligence depends on the value, timing, cost, and reliability of information across the complete system.

\section{Engineering Judgment}

The most valuable skill is not memorizing architectures. It is asking the correct questions:

\begin{quote}
What physical information is required? Which modality observes it? How is the evidence represented and aligned? Where is uncertainty handled? What fusion level is appropriate? What is the dominant resource cost? What happens when a sensor fails? Does the result reach the actuator before the deadline?
\end{quote}

These questions apply to autonomous vehicles, industrial inspection, rail systems, drones, medical devices, smart infrastructure, and edge sensing.

\begin{keyinsight}
A paper becomes part of a coherent technical world when it can be located by its physical inputs, information representation, fusion or estimation mechanism, hardware constraint, and effect on the final decision.
\end{keyinsight}

\section{Final Perspective}

Embedded AI is a systems discipline. Neural networks are important components, but the defining engineering challenge is the reliable conversion of uncertain physical evidence into timely action on constrained hardware. Building this complete view transforms literature review from collecting isolated algorithms into understanding a design space.
